\section{Summary and Future Work}

\subsection{Achieved Results}
The graduation project has successfully designed and implemented an AI-integrated Online Library Management System. The main achievements include:
\begin{itemize}
    \item \textbf{Modern Full-Stack Development:} Built a stable system comprising a Frontend (React.js, Vite), Backend (Node.js, Express, Prisma, PostgreSQL), and an AI Microservice (Python, FastAPI). The system is successfully deployed on cloud environments (Vercel, Google Cloud Run).
    \item \textbf{Effective AI Integration:} Successfully implemented Hybrid Semantic Search, an AI Virtual Assistant (Chatbot) utilizing RAG to answer library regulations, a personalized book recommendation system, and automated ISBN cataloging.
    \item \textbf{Comprehensive Library Digitization:} Completed core operational workflows including catalog management, physical copy inventory, borrow/return processes, renewals, reservations, inter-branch transfers, and multi-channel notification systems.
    \item \textbf{Optimization and Security:} Addressed fundamental security vulnerabilities such as IDOR, SQL Injection (via ORM), configured CORS, implemented Rate Limiting, and optimized database query performance (e.g., N+1 queries).
\end{itemize}

\subsection{Remarks and Evaluation}
\textbf{Strengths:}
\begin{itemize}
    \item The microservices-oriented architecture cleanly separates standard business logic (Backend) from intensive computational tasks (AI Service), facilitating scalability and ease of maintenance.
    \item User experience is significantly enhanced through AI features such as natural language search and 24/7 chatbot support.
    \item The system operates stably with an intuitive interface and fast response times, thanks to caching mechanisms and query optimizations.
\end{itemize}

\textbf{Limitations:}
\begin{itemize}
    \item Due to deployment on Free Tier infrastructure (Vercel and Google Cloud Run), the system occasionally experiences "Cold Starts" on initial requests or faces execution timeouts for complex AI tasks.
    \item The personalized recommendation algorithm has not yet reached its full potential due to the limited amount of historical borrowing and rating data in the testing environment (the Cold-Start problem for Recommender Systems).
    \item Some advanced features, such as automated book summarization, were developed but temporarily disabled due to API cost constraints and LLM rate limits.
\end{itemize}

\subsection{Future Work}
Based on the current limitations and expansion potential, the project proposes the following future development directions:
\begin{itemize}
    \item \textbf{Infrastructure Upgrade:} Migrate to dedicated paid cloud infrastructure (e.g., AWS or Google Cloud Platform) to eliminate Cold Starts, increase bandwidth, and support a larger concurrent user base.
    \item \textbf{AI Optimization and Expansion:} Research and fine-tune lightweight open-source LLMs (e.g., Llama 3 or Mistral) deployed internally to reduce reliance on third-party APIs (Gemini), thereby lowering costs and improving response times. Restore and upgrade the Function Calling capabilities of the Chatbot to allow users to query personal account information.
    \item \textbf{Enhanced Recommendation System:} Integrate advanced machine learning signals (e.g., Deep Learning-based recommender systems) and collect additional user behavior data (such as page view duration and click-through rates) to provide more accurate recommendations.
    \item \textbf{Mobile Application Development:} Build a cross-platform Mobile App using React Native or Flutter, integrating barcode/QR code scanning to assist librarians in inventory management and provide readers with faster on-the-go interactions.
\end{itemize}